\section{Constant-Time Oblivious Execution}
\label{sec:constant-time}

\subsection{Motivation and scope}

The normal profile executes only the blocks on the taken path. That is efficient,
but it makes the instruction trace and the private-memory access pattern depend on
secret inputs. For functions whose control flow itself is sensitive, Sirus offers
an optional fixed-pass oblivious mode selected by \verb|--constant-time|.

The mode targets a specific leakage class: the sequence of decoded blocks and the
sequence of private-memory addresses. It does not attempt to make every
microarchitectural effect input-independent. In particular, it does not model
data-dependent ALU latency, cache-bank conflicts, speculative execution, or
value-dependent power/EM leakage. It is not a formal Spectre mitigation.

\subsection{Frontier execution}

The manifest is parsed once. On every pass, every manifest block is decoded and
executed. The active path is carried in a frontier of $(block, mask)$ pairs. A
register write proposed by a block is applied as

\begin{equation}
  v_{\text{new}} = (v_{\text{proposed}} \land m) \lor (v_{\text{old}} \land \lnot m),
  \label{eq:predicated-write}
\end{equation}

where $m$ is the active mask (all ones or zero for a single path). A conditional
branch does not choose one successor; it enqueues both successors with masks
$m \land c$ and $m \land \lnot c$, where $c$ is the condition value. Returns and
traps are accumulated under the active mask. The loop is shown in
Algorithm~\ref{alg:ct}.

\begin{algorithm}[htbp]
\caption{Fixed-pass oblivious execution.}
\label{alg:ct}
\begin{algorithmic}[1]
\REQUIRE manifest $M$, entry state $e$, pass bound $P$
\ENSURE value and status
\STATE $F \gets \{(e,\; \text{all-ones})\}$; $R \gets 0$; $T \gets 0$
\FOR{$p = 1$ \TO $P$}
  \STATE $N \gets \emptyset$
  \FORALL{$(b,m) \in F$}
    \STATE decode and execute every instruction of block $b$ with mask $m$
    \STATE apply each register write with Eq.~\eqref{eq:predicated-write}
    \IF{block $b$ ends in a return with value $v$}
      \STATE $R \gets R \lor (m \land v)$
    \ELSIF{block $b$ ends in a trap}
      \STATE $T \gets T \lor m$
    \ELSIF{block $b$ ends in a conditional transition $(c,a,b',\text{mix})$}
      \STATE $N \gets N \cup \{(a,\; m \land c),\;(b',\; m \land \lnot c)\}$
    \ELSIF{block $b$ ends in an unconditional transition to $a$}
      \STATE $N \gets N \cup \{(a,\; m)\}$
    \ENDIF
  \ENDFOR
  \STATE $F \gets N$
  \IF{$F = \emptyset$}
    \STATE \textbf{break}
  \ENDIF
\ENDFOR
\IF{$F \neq \emptyset$}
  \STATE \RETURN $(0,\textsc{LoopBound})$
\ENDIF
\STATE \RETURN $(R,\textsc{Returned})$
\end{algorithmic}
\end{algorithm}

\subsection{Private-memory oblivious access}

A normal load or store uses the address in a virtual register. In constant-time
mode, the address is secret. The runtime therefore scans the entire private
memory for every byte of the access and selects the target with a mask:

\begin{equation}
  \text{selected} = \bigvee_{j=0}^{|\text{mem}|-1} \left( \text{mem}[j] \land
  [\,j = a + i\,] \right),
  \label{eq:oblivious-load}
\end{equation}

where $[\cdot]$ is the all-ones/all-zero mask of a predicate. A store writes back
every location: the selected location receives the new byte, and every other
location receives its old value. An inactive block uses address zero for its load
and suppresses its store, so a stale address cannot produce \code{BadPayload} or
corrupt memory. The access pattern is independent of the secret address; the cost
is proportional to the memory size, not to the number of bytes touched.

\subsection{Pass bound and convergence}

A program with no cycles has a finite maximum block count; the mode uses that
count as an exact bound. A program containing a loop cannot generally be bounded
without a loop limit, so the mode uses \verb|--constant-time-passes| (default
256, maximum 4096). The frontier is required to become empty within the bound. A
non-empty frontier after the last pass returns \code{LoopBound}. This is a
deliberate fail-closed behavior: the mode does not silently truncate a loop or
return a wrong value.

The bound is a compile-time constant in the generated runtime. It is visible to
an attacker, but it does not depend on the input. Functions with long or
data-dependent loop counts may require a larger bound or may not be suitable for
the mode.

\subsection{Property}

\begin{proposition}[Fixed instruction and address stream]
For a fixed protected program, fixed pass bound $P$, and fixed entry state, if
the frontier converges within $P$, then the sequence of decoded basic blocks and
the sequence of private-memory addresses in constant-time mode are independent of
the function's integer inputs. Register values, branch condition values, and the
number of passes actually executed before convergence may depend on the inputs.
\end{proposition}

\begin{proof}[Proof sketch]
The manifest order and the set of blocks are fixed at build time. The algorithm
executes every block in the frontier on every pass and never follows a single
branch edge; both successors are enqueued with complementary masks. The memory
access loop scans every byte position for every access. Therefore the decoded
block sequence and the address sequence are functions of the program and the pass
bound, not of the input values. Masked writes change values, not which
instructions or addresses are visited. If the frontier does not converge within
$P$, the runtime returns \code{LoopBound} rather than executing a variable number
of additional passes.
\end{proof}

The proposition is a source-level statement about the interpreter. It does not
claim that equal instruction and address streams imply equal execution time:
value-dependent ALU operations, cache behavior, branch prediction state, and
speculative execution are outside the model. The measured dudect and Cachegrind
results below are empirical checks of the source-level property and its practical
timing behavior on the test host.

\subsection{Measured behavior}

The release tree was measured with dudect-bencher and Cachegrind. On the final
recorded runs:

\begin{itemize}
  \item the cipher path measured \code{|t| = 2.39} after the two-key benchmark
        was changed to pass the key through \code{black_box}; without that
        control, the compiler specialized two constant-key paths and the
        benchmark measured code layout rather than cipher data timing;
  \item the four constant-time VM cases measured worst \code{|t| = 3.19}
        (straight-line 1.95, branch 1.67, loop 3.19, memory 2.73 in magnitude);
  \item the normal branchy VM profile measured \code{max |t| = 204}, confirming
        that the default profile is not constant-time;
  \item Cachegrind reported identical instruction references, data references,
        D1 misses, LLd misses, branches, and mispredicts for the branch-only
        constant-time leaf across two inputs.
\end{itemize}

The dudect threshold of \code{|t| = 5} is a conventional statistical screen, not
a proof. The measurements are host- and compiler-dependent and are reported with
their exact harness and seed.

\subsection{Cost}

The constant-time mode executes every block on every pass and scans the whole
private stack for each memory byte. On the benchmark leaf function it costs about
234 $\mu$s per call, compared with about 75 $\mu$s for the
default profile. The cost is expected: the mode trades data-dependent work for a
fixed work envelope. It is intended for small, high-value functions where the
timing property matters more than throughput.

\subsection{Limitations}

The mode is not a general constant-time compiler. It requires a finite pass
bound, a function that fits the private memory model, and the supported semantic
subset. It does not provide a proof of architectural timing equivalence, and it
does not address speculative execution or microarchitectural covert channels.
The normal profile remains the default for performance-sensitive deployments, and
its timing leak is documented rather than hidden.
